\chapter{Application to the Picard group} \label{XI}

This expos\'e is modelled on the preceding one, but this time the result of \numero \Ref{XI.1} is weaker.

Throughout this expos\'e, $X$ will denote a locally noetherian prescheme, $\Ii$ a quasi-coherent ideal of $\OX$, ($Y=V(\Ii)$ is therefore a closed subset of $X$), $U$ a variable open neighbourhood of $Y$ in $X$, and $\hat X$ the formal completion of $X$ along $Y$. For every ringed space $(Z, \Oo_Z)$, one denotes by $\PP(Z)$\refstepcounter{toto}\label{pXI.1} the category of invertible $\Oo_Z$-Modules, in other words locally free of rank~$1$, and by $\Pic(Z)$ the group of isomorphism classes of invertible Modules on $Z$.

\section{Comparison of $\Pic(\hat X)$ and of $\Pic(Y)$} \label{XI.1}

For every $n\in\NN$, set $X_n = (Y, \OX/\Ii^{n+1})$ and $P_n = \Ii^{n+1}/\Ii^{n+2}$. The sequence of sheaves of abelian groups on $Y$
\begin{equation} \label{eq:XI.1.1} 0\to P_n\lto{u}\Oo_{X_{n+1}}^*\lto{v}\Oo_{X_n}^*\to 1
\end{equation}
is \emph{exact}. Let us specify that the group structure of $P_n$ is the additive structure, that $u(x)=1+x$ for every $x\in P_n$, and that $v$ is the homomorphism deduced from the injection $\Ii^{n+2}\to\Ii^{n+1}$. One sees that $v$ is surjective by remarking that, for every $y \in Y$, $\Oo_{X_n, y}$ is a local ring, quotient of $\Oo_{X_{n+1}, y}$ by a nilpotent ideal; the rest is equally trivial. One deduces from (\Ref{eq:XI.1.1}) an exact cohomology sequence:
\begin{equation*} \label{eq:XI.1.*} \tag{$*$}
\H^1(Y, P_n)\lto{u^1}\H^1(Y, \Oo_{X_{n+1}}^*)\lto{v^1} \H^1(Y, \Oo_{X_n}^*)\lto{d}\H^2(Y, P_n){\quoi}.
\end{equation*}
Moreover, for every $n\in\NN$, one knows how to identify $\Pic(X_n)$ and $\H^1(Y, \Oo_{X_n}^*)$; moreover, if~$E$ is an invertible $\Oo_{X_{n+1}}$-Module, corresponding to a cohomology class $c(E)$, the cohomology class corresponding to the inverse image of $E$ on $X_n$ is equal to $v^1(c(E))$. Hence the following proposition:

\begin{proposition} \label{XI.1.1} Let us keep the notations introduced above. Let $p\in\NN$. The map $\Pic(\hat X)\to \Pic(Y_n)$:
\begin{enumeratei}
\item
is injective for $n\geq p$, if $\H^1(Y, P_n)=0$ for $n\geq p$;
\item
is an isomorphism for $n\geq p$, if $\H^i(Y, P_n)=0$ for $n\geq p$ and $i=1, 2$.
\end{enumeratei}
\end{proposition}

Of course, the exact sequence~\eqref{eq:XI.1.*} contains more information than the proposition above. The reader will have noticed that nothing has been said about the functor $\PP(\hat X)\to\PP(Y)$. Given two invertible $\Oo_{\hat X}$-Modules $E, F$, $H=\SheafHom(E, F)$ is also invertible. If one indicates by a subscript $n$ the reduction modulo $\Ii^{n+1}$, one finds an exact sequence:
\[
0\to H_0\otimes P_n \to \SheafHom(E_{n+1}, F_{n+1}) \to \SheafHom(E_n, F_n)\to 0{\quoi}.
\]
Hence an exact cohomology sequence which we shall not write and whose interpretation is evident; one may use this remark to study the functor~$\PP$.

\section{Comparison of $\Pic(X)$ and of $\Pic(\hat X)$} \label{XI.2}

The reader will find in expos\'e~\Ref{X}, \numero \Ref{X.2}, the proof of what follows:

\begin{proposition} \label{XI.2.1} Suppose that one has $\Lef(X, Y)$; then for every open neighbourhood $U$ of $Y$ in $X$, the functor
\begin{equation} \label{eq:XI.2.1} \PP(U)\to\PP(\hat X)
\end{equation}
is fully faithful, the map
\begin{equation} \label{eq:XI.2.2} \Pic(U)\to\Pic(\hat X)
\end{equation}
is therefore injective. If one has $\Leff(X, Y)$, the map~(\Ref{eq:XI.2.3}) is an isomorphism:
\begin{equation} \label{eq:XI.2.3}
\varinjlim_U\Pic(U)\to\Pic(\hat X){\quoi}.
\end{equation}
\end{proposition}

\begin{corollary} \label{XI.2.2} Suppose that one has $\Lef(X, Y)$ and that for every integer $n\geq p$, one has $\H^1(Y, P_n)=0$; then for every open $U\supset Y$, the maps
\[
\Pic(X)\to\Pic(U)\to\Pic(Y_n)
\]
are injective if $n\geq p$. If one has $\Leff(X, Y)$ and if moreover, for every integer $n\geq p$, one has $\H^i(Y, P_n)=0$ for $i=1$ and $i=2$, then the map
\[
\varinjlim_U\Pic(U)\to\Pic(Y_n)
\]
is an isomorphism for $n\geq p$.
\end{corollary}

\section{Comparison of $\PP(X)$ and of $\PP(U)$} \label{XI.3}

A definition:
\refstepcounter{subsection}\label{XI.3.1}
\begin{enonce*}[remark]{Definition 3.1\sfootnotemark}\sfootnotetext{For a more detailed study of the notion of parafactoriality, and the proof of~\Ref{XI.3.3}, \cf \EGA IV 21.13, 21.14.}
Let $X$ be a prescheme and let $Z$ be a closed subset of~$X$. Set $U=X-Z$. One says that $X$ is \emph{parafactorial} at the points of $Z$ if, for every open $V$ of~$X$, the functor $\PP(V)\to \PP(V\cap U)$ is an equivalence of categories. One also says that the pair $(X, Z)$ is parafactorial.
\end{enonce*}

Recall that $\PP(Z)$ denotes the category of Modules locally free of rank~$1$ on~$Z$.

\begin{definition} \label{XI.3.2}
A noetherian local ring is said to be \emph{parafactorial} if the pair $(\Spec(A), \{\rr(A)\})$ is parafactorial.
\end{definition}

One proves the following proposition which shows that the notion is ``pointwise'':

\begin{proposition} \label{XI.3.3}
Suppose $X$ locally noetherian. In order that the pair $(X, Z)$ be parafactorial, it is necessary and sufficient that, for every $z\in Z$, the local ring $\Oo_{X, z}$ be so.
\end{proposition}

Let us remark that in parafactorial there is ``fully faithful''. One proves as in lemma~\Ref{X.3.5} of expos\'e~\Ref{X} the:

\begin{lemma} \label{XI.3.4}
If $X$ is a locally noetherian prescheme and if $Z=X-U$ is a closed subset of $X$, the following conditions are equivalent:
\begin{enumeratei}
\item
for every open $V$ of $X$, the functor $\PP(V)\to \PP(V\cap U)$ is fully faithful;
\item
the homomorphism $\OX \to i_*(\Oo_U)$ is an isomorphism;
\item
for every $z\in Z$, one has $\prof(\Oo_{X, z})\geq 2$.
\end{enumeratei}
\end{lemma}

Thus ``parafactorial'' means that the conditions of~\Ref{XI.3.4} are satisfied and that, for every open $V$ of $X$, the homomorphism $\Pic(V)\to\Pic(V\cap U)$ is surjective. In particular, if $X$ is the spectrum of a noetherian local ring, one finds:

\begin{proposition} \label{XI.3.5}
Let $A$ be a noetherian local ring; in order that it be parafactorial it is necessary and sufficient that $\prof A\geq 2$ and $\Pic(X'-\{x\})=0$, where one has set $X'=\Spec(A)$ and where $x$ is the unique closed point of $X'$.
\end{proposition}

One will note that a local ring of dimension $\leq 1$ is never parafactorial since its depth is $\leq 1$. Thus \emph{factorial does not imply parafactorial}; it is however true for noetherian local rings of dimension $\geq 2$, as we shall see below.

\begin{lemma} \label{XI.3.6} Let $X$ be a locally noetherian prescheme and let $Z$ be a closed subset of $X$. Let $f\colon X_1\to X$ be a faithfully flat and quasi-compact morphism. Set $Z_1=f^{-1}(Z)$. If $(X_1, Z_1)$ is parafactorial, then $(X, Z)$ is so.
\end{lemma}

One remarks first that, if $i\colon(X-Z)\to X$ denotes the canonical immersion of $U=X-Z$ into $X$, the formation of the direct image by $i$ of a quasi-coherent $\Oo_U$-Module commutes with the base change $f$, since the latter is flat. It is therefore equivalent to assume the equivalent conditions of lemma~\Ref{XI.3.5} for $(X, Z)$ or for $(X_1, Z_1)$, since $f$ is a descent morphism for the category of quasi-coherent sheaves. It remains to prove that, for every open $V$ of $X$, $\Pic(V)\to\Pic(V\cap U)$ is surjective. One makes the base change $V\to X$, which changes nothing (\textit{sic}), and one is reduced to $V=X$. One then remarks that, if ${\Lcal}$ is an invertible $\Oo_U$-Module and if ${\Lcal}$ admits a locally free embedding, this extension is isomorphic to $i_*({\Lcal})$, by what has just been seen. It remains to prove that $i_*({\Lcal})$ is invertible. Using once more the fact that the direct image by $i$ commutes with flat base change, and that ``locally free of rank~$1$'' is a property which descends by a faithfully flat and quasi-compact morphism, we are done.

\begin{corollary} \label{XI.3.7} Let $A$ be a noetherian local ring; if $\hat A$ is parafactorial $A$ is so.
\end{corollary}

Do not believe that, if $A$ is parafactorial, $\hat A$ is so\nde{for a precise study of the relation between the factoriality of $A$ and of its completion, see (Heitmann~R., ``Characterization of completions of unique factorization domains'', \emph{Trans. Amer. Math. Soc.} \textbf{337} (1993), \numero 1, p\ptbl 379--387).}.

Before tackling the statement of the main theorem of this \no, let us make the connection with the theory of divisors and the notion of factorial ring\sfootnote{For the generalities which follow, \cf also \EGA IV~21.}.

Let $X$ be a \emph{noetherian} and \emph{normal} prescheme. Let $Z^1(X)$ be the free abelian group generated by the $x\in X$ such that $\dim\OXx=1$. The local ring of such a point is a discrete valuation ring. One will denote by $v_x$ the corresponding normalized valuation. Let $K(X)$ be the ring of rational functions of $X$ and let
\[
p\colon K(X)^*\to Z^1(X)
\]
be the map which to every $f\in K(X)^*$ associates the 1-codimensional cycle:
\[
(f) = \sum_{x\in X, \;\dim\OXx=1} v_x(f)\cdot x{\quoi}.
\]
The image of $p$ is denoted $P(X)$ and one calls its elements \emph{principal divisors}\sfootnote{In accordance with the terminology of \EGA IV~21, we now prefer to reserve the name of ``divisors'' for ``locally principal divisors'' or ``Cartier divisors''.}. One sets
\[
\Cl(X) = Z^1(X)/P(X){\quoi}.
\]
Let $Z^{\prime1}(X)$ be the subgroup of $Z^1(X)$ whose elements are the \emph{locally principal divisors}. One knows that
\[
\Pic(X)\simeq Z^{\prime1}(X)/P(X){\quoi},
\]
and consequently $\Pic(X)$ identifies with a subgroup of $\Cl(X)$.

Let us remark that if $U$ is a dense open of $X$, $K(X)\to K(U)$ is an isomorphism, and that if $\codim(X-U, X)\geq 2$, \ie if every $x\in X$ such that $\dim\OXx\leq 1$ belongs to $U$, the homomorphism $Z^1(X)\to Z^1(U)$ and consequently $\Cl(X)\to\Cl(U)$ is also an isomorphism. Finally, if every $x\in U$ is factorial, \ie $\OXx$ is so, then $Z^1(U) = Z^{\prime1}(U)$ hence $\Pic(U)\simeq\Cl(U)$.

\begin{propositionblah} \label{XI.3.7.1}
Let $X$ be a noetherian and normal prescheme. Let $(U_i)_{i\in I}$ be a family of opens of $X$ such that:
\begin{enumeratea}
\item
the $U_i$ form a filter base\nde{\ie a decreasing filtered family.};
\item
if one sets $Y_i=X-U_i$, one has $\codim(Y_i, X)\geq 2$ for every $i$,
\item
if $x\in U_i$ for every $i\in I$, then $\OXx$ is factorial.
\end{enumeratea}
\noindent
Then one has an isomorphism:
\[
\varinjlim_{i\in I}\Pic(U_i)\lto{\approx}\Cl(X){\quoi}.
\]
\end{propositionblah}
Let us remark that b)~implies that every $x\in X$ such that $\dim\OXx\leq 1$ belongs to $U_i$ for every $i$. Hence the $U_i$ are dense and moreover the homomorphism $Z^1(U_i)\to Z^1(X)$ is an isomorphism, as is $K(U_i)\to K(X)$. Hence $\Pic(U)\subset \Cl(U_i) \simeq \Cl(X)$. To prove what one wishes, it therefore suffices to show that every $D\in Z^1(X)$ belongs to $Z^{\prime1}(U_i)$ for a suitable $i$. It suffices to do this for the irreducible and positive ``divisors''. Let therefore $x\in X$ such that $\dim\OXx=1$. It suffices to prove that there exists an $i\in I$ such that $\overline{\{x\}}$ is locally principal at the points of $U_i$. Let $\Ii$ be the largest ideal of definition of the closed set $\overline{\{x\}}$. The set of points in a neighbourhood of which $\Ii$ is free is an open $U$. Now $U\supset \bigcap_{i\in I}U_i$ by~c). If one sets $Y=X-U$, one has $Y\subset \bigcup_{i\in I}Y_i$ with $Y_i=X-U_i$, now $Y$ is closed, hence admits a \emph{finite} number of generic points, hence is contained in the union of a finite number of $Y_i$, hence in $Y_j$ for some $j\in I$, since the $U_i$ form a filter base. Hence $U\supset U_j$.\qed

\begin{corollary} \label{XI.3.8}
Let $X$ be a noetherian and normal prescheme and let $Y$ be a closed subset of codimension $\geq 2$. Suppose that, for every $p\in X-Y$, $\Oo_{X, p}$ is factorial, then
\[
\Pic(X-Y)\to\Cl(X-Y)\to\Cl(X)
\]
are isomorphisms.
\end{corollary}

\begin{corollary} \label{XI.3.9}
Let $A$ be a noetherian and normal local ring. Set $X'=\Spec(A)$ and $x = \rr(A)$. In order that $A$ be \emph{factorial} it is necessary and sufficient that $\Pic(X'-\{x\})=0$ and that $\pp\in X'-\{x\}$ imply that $A_\pp$ is factorial.
\end{corollary}
Indeed, in order that $A$ be factorial it is necessary and sufficient that $\Cl(X')=0$\,\nde{see Bourbaki, \emph{Commutative algebra} VII.1.4, cor\ptbl of th\ptbl 2 and VII.3.2, th\ptbl 1}.

\begin{corollary} \label{XI.3.10}
Let $A$ be a noetherian local ring of dimension $\geq 2$. Set $X'=\Spec(A)$ and let $x = \rr(A)$. Set $X=X'-\{x\}$. The following conditions are equivalent:
\begin{enumeratei}
\item
$A$ is factorial;
\item
\textup{a)}\enspace for every $y\in X$, $\Oo_{X, p}$ is factorial, and
\par\hspace*{6.3mm}\textup{b)}\enspace $A$ is parafactorial, \ie $\prof A\geq 2$ and $\Pic(X)=0$.
\end{enumeratei}
\end{corollary}

Before proving this corollary, let us state the
\refstepcounter{subsection}\label{XI.3.11}
\begin{enonce*}{Serre's criterion for normality 3.11\sfootnotemark}\sfootnotetext{\Cf \EGA IV~5.8.6.} Let $A$ be a noetherian local ring. In order that $A$ be normal, it is necessary and sufficient that
\begin{enumeratei}
\item
for every prime ideal ${\pp}$ of $A$ such that $\dim A_{{\pp}}\leq 1$, $A_{{\pp}}$ be normal,
\item
for every prime ideal ${\pp}$ of $A$ such that $\dim A_{{\pp}}\geq 2$, one have $\prof A_{{\pp}}\geq 2$.
\end{enumeratei}
\end{enonce*}

Let us prove~\Ref{XI.3.10}.

(i)\ALORS (ii). Knowing that a localization of a factorial ring is again so, one has (ii)~a). Moreover $A$ is normal hence $\prof A\geq 2$ since $\dim A\geq 2$ (\Ref{XI.3.11}~(ii)). Finally $A$ is parafactorial; indeed $\Pic(X)\simeq\Cl(X')=0$ (\cf \Ref{XI.3.9}).

(ii)\ALORS (i). Let us first prove that $A$ is \emph{normal} by applying Serre's criterion. Since $\dim A\geq 2$, condition~(i) of the criterion is among the hypotheses. Moreover, for every ${\pp}\in X$, $A_{{\pp}}$ is factorial, hence normal, hence of depth $\geq 2$, at least if $\dim A_{{\pp}}\geq 2$. Finally $\prof A\geq 2$ by (ii)~b). It remains to apply~\Ref{XI.3.9}.

\smallskip
Let us summarize what precedes:

\begin{proposition} \label{XI.3.12}
Let $X$ be a locally noetherian prescheme and let $\Ii$ be a quasi-coherent ideal of $X$. Let $Y=V(\Ii)$. Let $p\in\NN$. Suppose that:
\begin{enumerate}
\item[1)] One has $\Leff(X, Y)$ (\textup{Expos\'e} \Ref{X});
\item[2)] $\H^i(X, \Ii^{n+1}/\Ii^{n+2})=0$ if $i=1$ or $2$ and if $n\geq p$;
\item[3)] for every open neighbourhood $U$ of $Y$ in $X$ and every $x\in X-U$, the ring $\OXx$ is parafactorial.
\end{enumerate}
\noindent
Then, for every $n\geq p$, and every open neighbourhood $U$ of $Y$, the homomorphisms
\[
\xymatrix{\Pic(X)\ar[rr]\ar[dr]&&\Pic(X_n)\\&\Pic(U)\ar[ur]&}
\]
are isomorphisms.
\end{proposition}

One knows some parafactorial rings:

\begin{theorem} \label{XI.3.13}
\begin{enumeratei}
\item
(Auslander-Buchsbaum)\nde{to be compared with the following purity result, due to Gabber. Let $X$ be the spectrum of a regular local ring $A$ of dimension $3$, $a$ of non-zero differential, \ie $a\in\mm-\mm^2$ and $U$ the complement of $V(a)$. Then a vector bundle on $U$ is free (for a simple proof, see (Swan~R.G., ``A simple proof of Gabber's theorem on projective modules over a localized local ring'' \emph{Proc. Amer. Math. Soc.} \textbf{103} (1988), \numero 4, p\ptbl 1025-1030). The rank~$1$ case is a particular case of theorem~\Ref{XI.3.13}. For purity results concerning vector bundles of arbitrary rank, whether in the analytic or algebraic setting, see (Gabber~O., ``On purity theorems for vector bundles'', \emph{Internat. Math. Res. Notices} (2002), \numero 15, p\ptbl 783--788).} A regular noetherian local ring is factorial (hence parafactorial if its dimension is $\geq 2$).
\item
A noetherian local ring of dimension $\geq 4$ which is a \emph{complete intersection} is parafactorial.
\end{enumeratei}
\end{theorem}

\begin{corollary}[Samuel's conjecture\label{XI.3.14}\ndemark]\ndetext{for a proof in the same vein, but more elementary, see Call~F.\ \&\ Lyubeznik~G., ``A simple proof of Grothendieck's theorem on the parafactoriality of local rings'', in \emph{Commutative algebra: syzygies, multiplicities, and birational algebra (South Hadley, MA, 1992)}, Contemp. Math., vol.~159, American Mathematical Society, Providence, RI, 1994, p\ptbl 15--18.} A noetherian local ring $A$ which is a complete intersection and which is factorial in codimension $\geq 3$ (\ie $\dim A_{{\pp}}\leq 3$ implies that $A_{{\pp}}$ is factorial) is factorial.
\end{corollary}

\emph{Let us prove the corollary}

One proceeds by induction on the dimension of $A$.

If $\dim A\leq 3$, $A$ is factorial by hypothesis.

If $\dim A>3$, by the induction hypothesis, remarking that a localization of a complete intersection is again so, all the localizations of $A$ other than $A$ are factorial. By theorem~\Ref{XI.3.13}~(ii), $A$ is parafactorial, hence factorial by~\Ref{XI.3.10}.

Let us prove~\Ref{XI.3.13}~(i) (following Kaplansky)\sfootnote{This is the proof reproduced in \EGA IV~21.11.1.}.

Let $A$ be a regular noetherian local ring, set $\dim A=n$. If $n=0$ or $1$, the result is known. Suppose $n\geq 2$, proceed by induction on $n$ and suppose $n\geq 2$ and the theorem proved for rings of dimension $<n$. Set $X'=\Spec(A)$ and $X=X'-\{x\}$, where $x=\rr(A)$. The localizations of $A$ other than $A$ are regular and of dimension $<n$, hence factorial. Moreover $\prof A = \dim A\geq 2$. It therefore suffices to prove that $\Pic(X)=0$ (cor\ptbl\Ref{XI.3.10}). Let therefore ${\Lcal}$ be an invertible $\OX$-Module, one knows that one can extend it to a coherent $\Oo_{X'}$-Module ${\Lcal}'$. There exists a resolution of ${\Lcal}'$ by free $\OX$-Modules:
\[
0\from {\Lcal}'\from {\Lcal}'_1\from\cdots\from{\Lcal}'_n\from 0{\quoi},
\]
since the cohomological dimension of $A$ is finite. By restriction to $X'$ one obtains a finite free resolution. It therefore suffices to prove the following lemma:

\begin{lemma} \label{XI.3.15} Let $(X, \OX)$ be a ringed space and let ${\Lcal}$ be a locally free $\OX$-Module which admits a finite resolution by modules \emph{free of finite type}. Then $\det({\Lcal}) \simeq \OX$.
\end{lemma}

Recall that one defines $\det({\Lcal})$ as the maximal exterior power of ${\Lcal}$. In the case under consideration, $\det({\Lcal}) \simeq{\Lcal}$ since ${\Lcal}$ is invertible, hence the lemma allows one to conclude. Let us prove this lemma. Let
\[
0\from {\Lcal}_0\from {\Lcal}_1\from {\Lcal}_2\from\cdots\from{\Lcal}_n\from 0{\quoi},
\]
be the announced exact sequence, where ${\Lcal}_0={\Lcal}$. Since everything is locally free, one has:
\[
\tbigotimes_{0\leq i\leq n}(\det({\Lcal}_i))^{(-1)^i} \simeq \OX
\]
now all the ${\Lcal}_i$, $i>0$, are free, hence also their determinants, hence also that of ${\Lcal}_0 = {\Lcal}$.\qed

\medskip
We still need to prove (ii) of the theorem. Beforehand, let us prove a lemma which will allow one to proceed by induction:

\begin{lemma} \label{XI.3.16}
Let $A$ be a noetherian local ring quotient of a regular ring. Let $t\in\rr(A)$ be an $A$-regular element. Suppose that $A$ is complete for the $t$-adic topology. Set $X'=\Spec(A)$, $Y'=V(t)\simeq \Spec(B)$, $B=A/tA$, $X=X'-\{x\}$, $Y=Y'-\{x\}$, $x=\rr(A)$. Suppose that:
\begin{enumeratea}
\item
for every $y\in X$ closed in $X$ one has $\prof\Oo_{X, y}\geq 2$,
\item
$\prof A/tA\geq 3$,
\end{enumeratea}
\noindent
then the map $\Pic(X)\to\Pic(Y)$ is injective. In particular, if $B$ is parafactorial, $A$ is so as well.
\end{lemma}

One knows that a)~implies $\Lef(X, Y)$ thanks to \Ref{X}~\Ref{X.2.1}. If one proves that $\H^1(Y, P_n)=0$ for every $n\geq 0$, one will know thanks to (\Ref{XI.2.2}) that $\Pic(X)\to\Pic(Y)$ is injective. If, moreover, $B$ is parafactorial, one will know that $\Pic(Y)=0$ (\Ref{XI.3.5}) hence $\Pic(X)=0$, and $\prof(A)=3+1\geq 2$ since $t$ is $A$-regular, hence $A$ will be parafactorial by~\Ref{XI.3.5}.

Let $\Ii=\widetilde{(tA)}$ be the $\Oo_{X'}$-Module associated with the ideal $tA$. In \numero \Ref{XI.1}, one set $P_n = (\Ii^{n+1}/\Ii^{n+2})_{|Y}$, for every $n\geq 0$. Now $t$ is $A$-regular hence $P_n\simeq \Oo_Y$. \emph{It therefore remains to prove that $\H^1(Y, \Oo_Y)=0$.} Now $Y = Y'-\{x\}$ is an open of $Y'$, one therefore has an exact sequence~(\Ref{I}~(\Ref{eq:I.27})):
\[
\H^1(Y', \Oo_{Y'})\to\H^1(Y, \Oo_{Y})\to\H^2_x(Y', \Oo_{Y'}){\quoi},
\]
whose right-hand term vanishes by virtue of hypothesis~b), and whose left-hand term vanishes because $Y'$ is affine.\qed

\begin{lemma} \label{XI.3.17}
Keeping the hypotheses of~\Ref{XI.3.16}, suppose moreover that:
\begin{itemize}
\item[c)] for every $y$ closed in $Y$, one has $\prof\Oo_{X, y}\geq 3$,
\item[d)] $\prof A/tA\geq 4$ (stronger than b),
\item[e)] for every closed $y$ of $X$, $y\in Y$, the ring $\Oo_{X, y}$ is parafactorial.
\end{itemize}
Then the map $\Pic(X)\to\Pic(Y)$ is an isomorphism, in particular in order that $A$ be parafactorial, it is necessary and sufficient that $B$ be so.
\end{lemma}

One knows~(\Ref{X}~\Ref{X.2.1}) that a)~and~c) imply $\Leff(X, Y)$. Moreover, by the reasoning which has just been given, d)~implies that $\H^1(Y, P_n)=0$ for every $n\geq 0$ and $i=1$ or $i=2$. Moreover, for every open neighbourhood $U$ of $Y$ in $X$, the complement of $U$ in $X$ consists of a finite number of closed points. Thanks to~e) and to theorem~\Ref{XI.3.12}, one deduces that $\Pic(X)\to \Pic(Y)$ is an isomorphism. On the other hand $\prof A\geq \prof B\geq 2$; by the criterion~\Ref{XI.3.5} one deduces that $A$ is parafactorial if and only if $B$ is so.

Let us now prove~\Ref{XI.3.13}~(ii). Let $R$ be a regular noetherian local ring. Let $(t_1, \ldots, t_k)$ be an $R$-sequence. Set $B=R/(t_1, \ldots, t_k)$ and suppose that $\dim B\geq 4$. One must prove that $B$ is parafactorial. One proceeds by induction on~$k$. If $k=0$, $B$ is regular, hence factorial by \Ref{XI.3.13}~(i), hence parafactorial by~\Ref{XI.3.10}. Suppose $k\geq 1$ and the theorem proved for $k'<k$. Set $A=R/(t_1, \ldots, t_{k-1})$ hence $B=A/t_k A$. One may assume $B$ complete by~\Ref{XI.3.7}. By the induction hypothesis, $A$ is parafactorial. Let us prove that one may apply lemma~\Ref{XI.3.17}. One has assumed $B$ complete hence also $A$, and therefore $A$ is complete for the $t$-adic topology. If $x\in X$, and if $x$ is closed in $X$, $A_x$ is a complete intersection of dimension $\geq 4$, with $k'<k$. By the induction hypothesis $A_x$ is parafactorial, and moreover of depth $\geq 4$. Which gives a), c)~and~e). Moreover, $\dim A\geq 5$, whence~d).\qed

\begin{theorem} \label{XI.3.18}
Let $X$ be a locally noetherian prescheme and let $\Ii$ be a coherent sheaf of ideals of $X$. Set $Y=V(\Ii)$. Let $n$ be an integer. Suppose that:
\begin{enumeratei}
\item
one has $\Leff(X, Y)$ (\cf examples \Ref{X}~\Ref{X.2.1}~and~\Ref{X}~\Ref{X.2.2}),
\item
for every $p\geq n$, one has $\H^i(Y, \Ii^{p+1}/\Ii^{p+2})=0$ for $i=1$ and $i=2$,
\item
for every open $U\supset Y$ and every $x\in X-U$, the ring $\OXx$ is regular of dimension $\geq 2$ or a complete intersection of dimension $\geq 4$.
\end{enumeratei}
\noindent
Then, for every open $U\supset Y$ and every integer $p\geq n$, the homomorphisms
\[
\Pic(X)\to\Pic(U)\to\Pic(Y_p)
\]
are isomorphisms, where $Y_p$ denotes the prescheme $(Y, \OX/\Ii^{p+1})$.
\end{theorem}

It suffices to combine \Ref{XI.3.12}~and~\Ref{XI.3.13}.
